%TC:incbib
\documentclass[11pt]{article} % This template uses article format. Change at your own risk
\usepackage{import}
\usepackage{authorinfo}
\usepackage{todonotes}
\usepackage{pdfpages}
\usepackage{tsetemplate} % This template is loaded as a package from tsetemplate.sty
% All of the Title Page and author info is defined in the referenced authorinfo file,
% if you want to change it, you should do it there, it can be found in the /definitions directory

\setcounter{secnumdepth}{-1}
%TC:fileinclude \import dir,file
%\import{definitions/}{authorinfo}

% Don't count words in citation macros. Citations themselves are not counted
%TC:macro \autocite 1
% ignore footnotes
%TC:macro \footnote 1

% This writes the word count to the word count file, so that the content of that file can be included (on the cover page).
\immediate\write18{texcount -inc -sum -1 main.tex > wordcount.tex}

% This is where the main document starts, if you need to put anything else in the preamble,
% then do it in tsetemplate.sty
\begin{document}

% We use a custom title page that fills the full first page and has no number but includes TSE logo.
% You must fill in the data in definitions/authorinfo.tex to fill these parameters
\import{definitions/}{mytitlepage}

% Start main document area, this will use the files from /section files  you can use include pages, you can ignore sections by adding
% NoSection to the relevant file, or just comment it out here

\singlespacing % why this is needed I don't know.
\renewcommand\contentsname{Table of Contents}
\tableofcontents{}
\doublespacing
\newpage

% \begin{centering}
% \noindent
% \color{section-label-color}
% \textbf{\huge Commentary}
% %\vspace{\baselineskip}
% \end{centering}

% Based on whether you have defined them in authorinfo.tex, your section files will be included or excluded here.
% To make this main document clearer, include things like tables or figures that have complex formatting into separate files in
% Subdirectory - here section files (don't add the filename extension eg .tex as a general good practice as it only works
% with import or input, not with include). - you can then include them with \import{/sectionfiles}{nameoffile}
% You can write directly into here, though it may be better to use the defined sub-files - you can use more and import or include them
% here - if you use include, put them in the same directory as this file, it's just like typing here.
% Use import if you have them in a subdirectory.
% In a big project you could include each chapter and its images and figures in it's own directory and then import them all
% you can use sub-import to include the other parts into the sub-file
% As with the TSE template, each section is a main unnumbered heading
% if you want numbered sections take out the * from \section  and comment the \addcontentsline
% or you will have duplicate entries in your Table of Contents

\ifdefined\Introduction
\phantomsection
\section*{\Introduction}
\addcontentsline{toc}{section}{\Introduction}
\label{intro}
\import{sectionfiles/}{introduction}
\fi

\ifdefined\CompA
  \phantomsection
  \section*{\CompA}
  \addcontentsline{toc}{section}{\CompA}
  \label{comp:a}
  \import{sectionfiles/}{alicia}
\fi

\ifdefined\CompB
  \phantomsection
  \section*{\CompB}
  \addcontentsline{toc}{section}{\CompB}
  \label{comp:b}
  \import{sectionfiles/}{gestrals}
\fi

\ifdefined\CompC
  \phantomsection
  \section*{\CompC}
  \addcontentsline{toc}{section}{\CompC}
  \label{comp:c}
  \import{sectionfiles/}{une-vie}
\fi

\ifdefined\Conclusion
\phantomsection
\section*{\Conclusion}
\addcontentsline{toc}{section}{\Conclusion}
\label{conc}
\import{sectionfiles/}{conclusion}
\fi

% \ifdefined\ResearchSection
%   \phantomsection
%   \section*{\ResearchSection}
%   \addcontentsline{toc}{section}{\ResearchSection}
%   \label{research}
%   \import{sectionfiles/}{researchsection}
% \fi


% \ifdefined\CreativeApproach
%   \phantomsection
%   \section*{\CreativeApproach}
%   \addcontentsline{toc}{section}{\CreativeApproach}
%   \label{creative}
%   \import{sectionfiles/}{creativesection}
% \fi

% \ifdefined\TechApproach
%   \phantomsection
%   \section*{\TechApproach}
%   \addcontentsline{toc}{section}{\TechApproach}
%   \label{tech}
%   \import{sectionfiles/}{techsection}
% \fi

% \ifdefined\PracticalApproach
%   \phantomsection
%   \section*{\PracticalApproach}
%   \addcontentsline{toc}{section}{\PracticalApproach}
%   \label{practical}
%   \import{sectionfiles/}{practicalsection}
% \fi

% \ifdefined\ProjectContext
%   \phantomsection
%   \section*{\ProjectContext}
%   \addcontentsline{toc}{section}{\ProjectContext}
%   \label{context}
%   \import{sectionfiles/}{projectcontextsection}
% \fi

% \ifdefined\FurtherComments
%   \phantomsection
%   \section*{\FurtherComments}
%   \addcontentsline{toc}{section}{\FurtherComments}
%   \label{further}
%   \import{sectionfiles/}{furthercommentsection}
% \fi

%TC:ignore (don't include the word count statement in the actual word count)
\textit{\Wordcount words.}
%TC:endignore

\clearpage

%TC:ignore
% The Bibliography will be added at this point.

\phantomsection
\printbibliography[title={Bibliography}, heading=bibintoc] %
\label{references}
% If you only want sources you've cited in the text, then comment out the below \nocite
% otherwise, it will cite any and all sources in your citation file.
% this might not be very useful if you are using it for multiple projects
% \nocite{*}


% \ifdefined\Endmatter
%   \clearpage % This will put this section at the start of a new page.
%   \phantomsection
%   \section*{\Endmatter}
%   \addcontentsline{toc}{section}{\Endmatter}
%   \label{endmatter}
%   \singlespacing
%   \import{sectionfiles/}{endmattersection}
% \fi

\ifdefined\MyAppendixA
  \clearpage % This will put this section at the start of a new page.
  \phantomsection
  \section*{\MyAppendixA}
  \addcontentsline{toc}{section}{\MyAppendixA}
  \label{appendixA}
  \singlespacing
  \import{sectionfiles/}{appendix-a}
\fi

\clearpage

\ifdefined\MyAppendixB
  \clearpage % This will put this section at the start of a new page.
  \phantomsection
  \section*{\MyAppendixB}
  \addcontentsline{toc}{section}{\MyAppendixB}
  \label{appendixB}
  \doublespacing
  \import{sectionfiles/}{appendix-b}
\fi

\clearpage

\ifdefined\MyAppendixC
  \clearpage % This will put this section at the start of a new page.
  \phantomsection
  \section*{\MyAppendixC}
  \addcontentsline{toc}{section}{\MyAppendixC}
  \label{appendixC}
  \singlespacing
  \import{sectionfiles/}{appendix-c}
\fi
\clearpage

\ifdefined\MyAppendixD
  \clearpage % This will put this section at the start of a new page.
  \phantomsection
  \section*{\MyAppendixD}
  \addcontentsline{toc}{section}{\MyAppendixD}
  \label{appendixD}
  \singlespacing
  \import{sectionfiles/}{appendix-d}
\fi
\clearpage

%%%% APPENDICIES %%%%
%% Copy/Paste and uncomment, and replace N with whatever you've defined in authorinfo.tex
% Make a filename that matches called appendix-n (again replace n with the correct letter)
% in the files section

% \ifdefined\MyAppendixN
%    \clearpage % This will put this section at the start of a new page.
%   \addcontentsline{toc}{section}{\MyAppendixN}
%   \section*{\MyAppendixN}
%   \singlespacing
%   \import{sectionfiles/}{appendix-n}
% \fi

%%%% List of Figures and tables.
%%% IF you want to include a list of figures and / or tables, you can do that below.
%%% Comment it out as appropriate.

\renewcommand{\listfigurename}{List of Figures}
\renewcommand{\listtablename}{List of Tables}
\cleardoublepage
% \phantomsection
% \addcontentsline{toc}{section}{\listfigurename}
% \label{listfigures}
% \listoffigures
% \phantomsection
% \addcontentsline{toc}{section}{\listtablename}
% \label{listtables}
% \listoftables
\doublespacing


%TC:endignore
\end{document}
